\section{Experimental Setups}
\label{sec:experiments}

\begin{figure*}[t]
\centering
\includegraphics[width=0.9\linewidth]{latex/figures/new_main_result_tagged.pdf}
\caption{\crchange{Quality--latency trade-off on ACL 60/60. Terminology accuracy uses strict exact-form matching. \method improves terminology accuracy in all three directions; BLEU changes are positive in most, but not all, settings.}}
\label{fig:main_result_1}
\vspace{-0.5cm}
\end{figure*}

\subsection{Data}
\label{sec:task_datasets}

\paragraph{Retriever Training} We synthesize the retriever training data from GigaSpeech~\citep{chen21o_interspeech} and Wikidata\footnote{\url{https://www.wikidata.org/}}. GigaSpeech is a multi-domain English ASR corpus containing 10K hours of high-quality labeled audio. We process it following the procedure in Section~\ref{sec:retriever}, yielding 4.7M training pairs. The variable speech-context duration is sampled uniformly from ${2.88, 3.84, 4.80, 5.76}$ seconds, corresponding to the inference-time chunk sizes ${0.96, 1.92, 2.88, 3.84}$ plus a $t_\text{lb}=1.92$-second look-back. The look-back captures terms that cross chunk boundaries and we set it to $1.92$ seconds since most terms are shorter than this, as shown in Figure~\ref{fig:term_duration}.

For Wikidata, we extract candidate terms from the RDF truthy dump, keeping entities that have both an English \texttt{rdfs:label} and at least one \texttt{wdt:P31} triple. Full processing details are given in Appendix~\ref{sec:apdx_wikidata_proc}. For each retained term, we prompt Gemini-2.5-Flash to generate three natural utterances of 5--15 words containing the term (Figure \ref{fig:wikidata_utt_gen}). We then synthesize the utterances with CosyVoice 3\footnote{FunAudioLLM/Fun-CosyVoice3-0.5B-2512}, using speaker prompts randomly sampled from GigaSpeech. This results in 2.9M training pairs.

\paragraph{\crchange{Speech} LLM Training}
We synthesize the Speech LLM training data using the same GigaSpeech ASR dataset. The interleaved En-Zh/De/Ja translation data is synthesized similar to InfiniSST with details in Appendix \ref{sec:speechllm_data_synth}. 
We subsample 12.5K instances for Speech LLM fine-tuning. We build a 140K glossary from GigaSpeech and Wikidata for distractor retrieval, and cap the size of $G_i$ at 20 terms.

\paragraph{Evaluation} \begin{crblock}We evaluate on two datasets. The ACL 60/60 dev set~\citep{salesky-etal-2023-evaluating} contains five English ACL talks translated into 10 target languages; we evaluate En--Zh/De/Ja and feed each unsegmented talk directly to the model. Its terminology lists are non-exhaustive. Source spans were automatically tagged, English technical terms were checked by domain experts, and target translations were post-edited by native professional translators and verified by a second annotator. We use these human-curated entries as the retrieval glossary; after source-side deduplication it contains 238 terms.

The ESO test set~\citep{eso_dataset} comprises five long-form English oncology training videos with human translations into French, Spanish, German, and Slovene. We report only En--De because the dataset does not provide human references for Chinese or Japanese. We use GPT-5.4 to extract and filter domain-specific glossary candidates---including procedures, diagnoses, biomarkers, clinical-trial concepts, and acronyms---and manually verify the retained entries. The final glossary has 215 rows corresponding to 212 unique source terms; Appendix~\ref{app:eso_glossary} gives the procedure and prompts. To test larger retrieval banks, we expand the ACL and ESO glossaries to 1K and 10K entries using unseen Wikidata terms from related domains.\end{crblock}

\subsection{Evaluation Metrics}
\label{sec:evaluation}

\paragraph{Latency}
Following IWSLT 2025 practice~\citep{agostinelli-etal-2025-findings}, we use StreamLAAL~\citep{papi-etal-2024-streamatt} as our latency metric. StreamLAAL segments the model hypothesis to align with reference sentences and reports the average latency over the resulting (hypothesis, reference) pairs.

\paragraph{Translation Quality}
\begin{crblock}
We evaluate aligned (hypothesis, reference) sentence pairs with SacreBLEU~\citep{papineni-etal-2002-bleu,post-2018-call} and XCOMET-XXL~\citep{guerreiro-etal-2024-xcomet}. We report XCOMET multiplied by 100. BLEU measures surface-form overlap, whereas the reference-based XCOMET model additionally assesses context-sensitive translation errors.
\end{crblock}

\paragraph{Terminology Accuracy}
\begin{crblock}
Our primary metric is exact-form terminology accuracy (TERM\_ACC). A glossary entry is eligible in an aligned sentence when its source form occurs in the source and its prescribed translation occurs in the reference. After deduplicating identical target forms, each eligible target form is counted at most once per sentence; TERM\_ACC is the percentage that appear verbatim in the hypothesis.

Because strict matching can reject valid orthographic or inflectional variants, we also report a conservative form-aware diagnostic. It accepts Unicode width, case, spacing, and hyphenation variants in all languages and contiguous German lemma sequences produced by a fixed morphological analyzer. It does not accept synonyms or semantic paraphrases. Appendix~\ref{app:form_aware_metric} provides the complete protocol and results.
\end{crblock}



\subsection{Model Configuration}
\label{sec:retriever_training}

\paragraph{Architecture}
For the retriever, we use a Qwen3-Omni Audio Transformer\footnote{\url{https://huggingface.co/Atotti/Qwen3-Omni-AudioTransformer}} as the speech encoder, followed by a linear projection. For text encoding, we use the dense BGE-M3\footnote{\url{https://huggingface.co/BAAI/bge-m3}} encoder and \crchange{use the contextualized representation at its first \textsc{[CLS]} position as the term vector}. For the Speech LLM, we use Qwen3-Omni-30B-A3B-Instruct\footnote{\url{https://huggingface.co/Qwen/Qwen3-Omni-30B-A3B-Instruct}}.

\paragraph{Training}
We train both the cross-modal retriever and the Speech LLM with LoRA~\citep{hu2022lora}. For retriever training, the number of hard negative $N$ is 1024 and the global batch size is 8192. For Speech LLM training, we use sequence packing and a global batch size of 8192 tokens. Full hyperparameters and details are provided in Appendix~\ref{app:training_details}.
% Dev-only threshold-selection evidence is available at latex/figures/retriever_dev_pr_fixedraw_hn_comparison.pdf.

\paragraph{Inference}

We implement the inference agent in SimulEval~\citep{ma-etal-2020-simuleval}, varying the chunk size $l$ over ${0.96, 1.92, 2.88, 3.84}$ seconds. The Speech LLM is invoked once per chunk and generates up to 40 new tokens per call. We serve it with vLLM~\citep{kwon2023efficient} using tensor parallelism of size 2, decode with temperature $0.6$, top-$p{=}0.95$, and top-$k{=}20$, and strip the assistant's \texttt{<term>} tags from the hypothesis before scoring.

For retrieval, we build a FAISS inner-product index over $\ell_2$-normalized term embeddings. At each step, the retriever encodes the current speech chunk together with a $1.92$-second look-back, and queries the index using only those windows whose right boundary lies within the current chunk. We use Top-$K{=}10$ and discard candidates with similarity below $\tau{=}0.78$\footnote{We select $\tau$ on the dev set to allow at most a 1\% drop in recall.}.


\begin{figure}[t]
\centering
\includegraphics[width=0.8\linewidth]{latex/figures/glossary_bank_ablation_zh_fixedraw.pdf}
\caption{Scaling the glossary from 0.2K to 1K and 10K terms. \method shows only a slight drop in terminology accuracy and translation quality, remaining well above the no-retrieval baseline.}
\label{fig:ablation_glossary_bank}
\vspace{-0.5cm}
\end{figure}


% \begin{figure}[t]
%   \centering
%   \includegraphics[width=\linewidth]{latex/figures/rag_compute_rtf.pdf}
%   \caption{RTF}
%   %\caption{RAG compute real-time factor under timeline retrieval. We compute $\mathrm{RTF}_{\mathrm{RAG}} = t_{\mathrm{retrieve}} / (0.96\,\mathrm{s}\times\mathrm{LM})$ from verified En-Zh medicine hardraw runtime traces. The retriever runs once per vLLM generation step and encodes the current generation span plus a fixed $1.92\,\mathrm{s}$ look-back. Bars show median per-call RTF, the magenta line shows mean RTF, and the gray line shows median retriever time. Even as the encoded span grows from $2.88$ to $5.76$ seconds, the median RAG compute share falls from $3.85\%$ to $1.14\%$.}
%   \label{fig:rag_compute_rtf}
% \end{figure}

\subsection{Baselines}

\paragraph{Offline}
We use Qwen3-Omni-30B-A3B-Instruct to translate pre-segmented speech utterances offline, with the same decoding parameters as \method. This serves as an upper bound on translation quality.

\paragraph{Offline + GT Terms}
We extend the \textbf{Offline} baseline by adding the ground-truth terms appearing in each utterance to the prompt in a manner similar to \citet{gong24b_interspeech}, helping the model handle terminology correctly. This serves as an upper bound on terminology translation accuracy.

\paragraph{InfiniSST}
InfiniSST~\citep{ouyang-etal-2025-infinisst} is a strong SST baseline for simultaneous translation of unbounded speech and achieved top performance in the low-latency track at IWSLT 2025~\citep{ouyang-etal-2025-cmus,agostinelli-etal-2025-findings}. We train InfiniSST on the same data and with the same Speech LLM as \method, but remove all terminology input $G_i$. This setting can be viewed as \method without retrieval.
